\section{PCD}


\begin{frame}[fragile]{Implementation of the polar coordinates decoupling (PCD). To polar}
  \begin{columns}[T,onlytextwidth]
    \column{0.55\textwidth}
      \begin{adjustbox}{max height=0.65\textheight,keepaspectratio}
\begin{lstlisting}[language=Python, basicstyle=\ttfamily\scriptsize, aboveskip=2pt, belowskip=2pt, numbers=left, numberstyle=\tiny, frame=single, breaklines=true]
def to_polar(x: torch.Tensor) -> Tuple[torch.Tensor, torch.Tensor]:
    p, k = x.shape
    magnitudes = torch.linalg.vector_norm(x, dim=1, keepdim=True)
    phis = torch.empty(p, k-1, dtype=x.dtype, device=x.device)
    x_squares = x * x
    x_squares_flipped = torch.flip(x_squares, dims=[1])
    x_squares_cumsum_flipped = torch.cumsum(x_squares_flipped, dim=1)
    x_squares_cumsum = torch.flip(x_squares_cumsum_flipped, dims=[1])
    y_head = x_squares_cumsum[:, 1:-1].sqrt()
    x_head = x[:, :-2]
    phis[:, :-1] = torch.atan2(y_head, x_head)
    phi_last = torch.atan2(x[:, -1], x[:, -2])
    phi_last = (phi_last + 2 * math.pi) % (2 * math.pi)
    phis[:, -1] = phi_last
    return phis, magnitudes.squeeze(-1)
\end{lstlisting}
      \end{adjustbox}
    \column{0.4\textwidth}
        The following transformation to spherical coordinates was implemented\\
        This function converts the vector representation $x$ from Cartesian coordinates to polar coordinates: it calculates the norm (modulus) of each vector and the angles between the axes using arctangents. For the last two coordinates, the angle is additionally adjusted so that it always lies in the range $(0,2\pi]$
  \end{columns}
\end{frame}
